\documentclass[11pt]{article} % Default font size
\pagenumbering{gobble}
\usepackage[skip=4pt plus1pt]{parskip}

\usepackage[sorting=ydnt,maxbibnames=99]{biblatex}
\newcommand{\myname}{\underline{Minsung Cho}}

\nocite{*}
\addbibresource{pubs.bib}

\usepackage[hmargin=20mm, top=14mm, bottom=12mm]{geometry} % US letter; set margins

\usepackage{hyperref}

\usepackage{titlesec}
\titlespacing*{\section}{0pt}{*1.0}{*1.3}

\setcounter{secnumdepth}{0} % Suppress section numbering

\newcommand{\tabbedblock}[1]{
\begin{tabbing}
\hspace{2.5cm} \= \hspace{4cm} \= \kill
#1
\end{tabbing}
}
%----------------------------------------------------------------------------------------

\begin{document}

%----------------------------------------------------------------------------------------
%	NAME AND CONTACT INFORMATION
%----------------------------------------------------------------------------------------

\parbox{0.5\textwidth}{ % First block
\huge{\textbf{Minsung Cho}}
}
\hfill % Horizontal space between the two blocks
\parbox{0.5\textwidth}{ % Second block
\begin{tabbing} % Enables tabbing
\hspace{3cm} \= \hspace{4cm} \= \kill % Spacing within the block
\href{mailto:minsungc1@gmail.com}{\texttt{minsungc1@gmail.com}} \\ % Home phone
\href{https://minsungc.github.io}{minsungc.github.io} \\
\href{https://linkedin.com/in/minsung-c}{linkedin.com/in/minsung-c}
\end{tabbing}}

% \vspace{2pt}
% \noindent Software engineer working on query languages, compilers, and database engines.

\hrule

%----------------------------------------------------------------------------------------
%	EMPLOYMENT HISTORY SECTION
%----------------------------------------------------------------------------------------

\newcommand{\job}[7]{%
{\setlength{\parskip}{0.2pt}%
\par\noindent\textbf{\textit{#6}}\hfill\textbf{#1 #2}\par
\noindent\href{#5}{#3}\hfill\textit{#4}\par}
\nopagebreak
\vspace{-3pt}
\begin{itemize}\setlength{\itemsep}{0.2pt}\setlength{\parskip}{0pt}\setlength{\topsep}{0pt}\setlength{\parsep}{0pt}
    #7
\end{itemize}
\vspace{0.3em}
}

\section{Employment}

\job
{May 2025 --}{Present}
{RelationalAI, Logic Engine team}{Remote}
{http://relational.ai}
{Software Engineer, Compiler}
{
    \item Design and implement the procedural extension to RelationalAI's query language
    \begin{itemize}
        \item Scale adoption from internal algorithms to several enterprise workloads
        \item Optimize previously intractable queries to sub-hour runtimes
    \end{itemize}
    \item Build the GQL-style graph pattern matching extension for the query language
    \item Maintain the database engine's Protocol Buffer-based entry point ($\sim$100k queries/day)
    \item Develop an event-based user-facing logging tool to diagnose slow queries
}

\job
{Sep 2022 --}{Apr 2025}
{Northeastern University, Programming Research Lab}{Boston, USA}
{http://prl.khoury.northeastern.edu}
{Graduate Researcher}
{
    \item Designed a language to express optimization problems under uncertainty~\cite{cho2025}
    \begin{itemize}
        \item Implemented a novel binary decision diagram algorithm to compute optimal values
        \item Outperformed a state-of-the-art probabilistic programming language by up to 10x
    \end{itemize}
    \item Co-led the first empirical study of online proof assistant communities~\cite{lincroft2024}
    \begin{itemize}
        \item Applied established developer lifecycle models to characterize contribution patterns
    \end{itemize}
}

\job
{May 2024 --}{Dec 2024}
{RelationalAI, Logic Engine team}{Remote}
{http://relational.ai}
{R\&D Intern}
{
    \item Shipped a decompiler that lifts optimized Datalog back to source query
    \item Integrated it into the production query profiler to automatically flag performance bottlenecks
}

% \job
% {Jun 2021 --}{Jul 2021}
% {University of Tennessee}{Chattanooga, USA}
% {http://www.utc.edu}
% {REU Researcher}
% {
%     \item Published research in the theory of matrices in function spaces~\cite{cho2021krein}
%     \item Automated routine computation in Mathematica, saving weeks of manual effort
% }

%----------------------------------------------------------------------------------------
%	EDUCATION SECTION
%----------------------------------------------------------------------------------------

\newcommand{\education}[6]{%
{\setlength{\parskip}{0pt}%
\par\noindent\textbf{\textit{#6}}\hfill\textbf{#1 #2}\par
\noindent\href{#5}{#3}\hfill\textit{#4}\par}
\nopagebreak
\vspace{0.3em}
}

\section{Education}

\education
{Sep 2022 -}{May 2025}
{Northeastern University}{Boston, USA}
{http://www.northeastern.edu}
{Master of Science, Computer Science}

\education
{Aug 2018 -}{May 2022}
{Carnegie Mellon University}{Pittsburgh, USA}
{http://www.cmu.edu}
{Bachelor of Science, Mathematics and Philosophy}

%----------------------------------------------------------------------------------------
%	Skills Section
%----------------------------------------------------------------------------------------

\section{Skills}

\begin{description}\setlength{\itemsep}{2pt}\setlength{\parskip}{0pt}\setlength{\topsep}{0pt}
    \item[Programming languages:] Julia, Python, Datalog, OCaml, Rust, Go, C/C++, Haskell, Standard ML
    \item[Tools:] Unix, Git, Jira, Docker, Datadog, Observe, Snowflake, LLVM/MLIR
\end{description}

%----------------------------------------------------------------------------------------
%	Awards Section
%----------------------------------------------------------------------------------------

\section{Awards and Recognition}

\noindent First Place, ACM Student Research Competition at PLDI (top PL conference)\hfill\textbf{2023}\par
\noindent Scholarship Recipient, Programming Languages Mentoring Workshop (PLMW) at PLDI\hfill\textbf{2023}\par
\noindent University and College Honors, Carnegie Mellon University\hfill\textbf{2022}\par

%----------------------------------------------------------------------------------------
%	Publication Section
%----------------------------------------------------------------------------------------
\section{Publications}

{\small
\setlength{\bibitemsep}{1pt}
\printbibliography[heading=none]}

\end{document}